\subsection{序向量空间上的拓扑}

设 E 是序向量空间。E 上的拓扑称为与 E 的序向量空间结构相容，如果它既与 E 的向量空间结构相容，又满足下列公理：

(TO) 诸 \(x\)（\(x \geqslant 0\)）组成的凸锥在 E 中是闭的。

赋有相容拓扑的序向量空间 E 称为序拓扑向量空间。

例. --- 赋有通常拓扑及其因子的序结构之乘积这一序结构的空间 \(R^{n}\) 是序拓扑向量空间。另一方面，当 \(n \geqslant 2\) 且 \(R^{n}\) 赋有字典序 (S, III, § 2.6) 时，通常拓扑与 \(R^{n}\) 的序向量空间结构不相容。

设 \(\mathbf{A}\) 是一个集；向量空间 \(\mathcal{B}(\mathbf{A};\mathbf{R})\)（由定义在 \(\mathbf{A}\) 上的实值有界函数组成），赋有由范数 \(\| x\| = \sup_{t\in \mathbf{A}}|x(t)|\) 定义的拓扑以及由

\(R^{A}\) 的乘积序结构诱导的序结构，是一个序拓扑向量空间。

在序拓扑向量空间 E 中，由元素 \(x \leqslant 0\) 组成的集是闭的；由于平移是同胚，我们推出，对所有 \(a \in E\)，由元素 \(x \geqslant a\)（相应地 \(x \leqslant a\)）组成的集是闭的。由于 \(\{0\}\) 是集 \(x \geqslant 0\) 与 \(x \leqslant 0\) 之交，由此推出 \(\{0\}\) 是闭的，从而 E 是 Hausdorff 的。

命题 18. --- 在序拓扑向量空间 E 中，设 H 是由关系 \(\leqslant\) 定向的一个集。若 H 的截段滤子在 E 中有极限，则该极限是 H 的上确界。

因为，设 \(b = \lim_{x\in \mathbf{H}}x\)；对每个 \(y\in \mathbf{H}\)，诸 \(x\in \mathbf{H}\)（满足 \(x\geqslant y\)）组成的集是

H 的截段滤子中的一个集，因此 b 是该集的一个接触点；但由于集 \(x \geqslant y\) 在 E 中是闭的，我们有 \(b \geqslant y\)，故 b 是 H 的一个上界。另一方面，若 a 是 H 的一个上界，则 H 包含在闭集 \(x \leqslant a\) 中；由于 b 是 H 的一个接触点，我们有 \(b \leqslant a\)，证毕（II, 第 72 页, exerc. 42）。

